% 図1：Salmon Ladder の概要（1段幅 8.4 cm 以内）．(a) 1つのグラフと支持画像からの質量 u+，(b) 参照画像での順位 p+ への BH 法で種を選ぶ（プール全体，折なし），(c) 種からの質量 u- と参照画像による読み出し sigma_ref, p_fin．
\providecommand{\figfont}{\fontsize{6.74pt}{8pt}\selectfont}
\begin{tikzpicture}[font=\figfont, >={Latex[length=1.3mm]}, x=1cm, y=1cm,
  sup/.style={draw=black, fill=black!75, minimum size=0.21cm, inner sep=0pt},
  idsq/.style={draw=black!55, fill=black!10, minimum size=0.19cm, inner sep=0pt},
  oodA/.style={draw=blue!60!black, fill=blue!35, minimum size=0.22cm, inner sep=0pt, regular polygon, regular polygon sides=3},
  seed/.style={draw=red!70!black, fill=red!45, minimum size=0.22cm, inner sep=0pt, regular polygon, regular polygon sides=3},
  calp/.style={draw=blue!70!black, fill=white, line width=0.6pt, minimum size=0.17cm, inner sep=0pt, circle},
  refp/.style={draw=teal!70!black, fill=teal!30, line width=0.5pt, minimum size=0.2cm, inner sep=0pt, diamond},
  edge/.style={black!45, line width=0.45pt},
  flow/.style={->, orange!85!black, line width=0.7pt},
  nflow/.style={->, red!70!black, line width=0.7pt},
  box/.style={draw, rounded corners=2pt, align=center, inner sep=2pt},
  lab/.style={inner sep=1pt, font=\figfont}]
% ================================================================ (a) one graph, mass from the support
\node[lab, anchor=west, font=\figfont\bfseries] at (0,6.6) {(a) 1つのグラフ：支持画像の質量 $u^{+}$ が辺を伝わる};
\node[sup] (s1) at (0.55,5.8) {}; \node[sup] (s2) at (1.05,6.05) {}; \node[sup] (s3) at (1.45,5.65) {};
\node[calp] (c1) at (0.9,5.35) {}; \node[refp] (c2) at (1.9,5.2) {};
\node[idsq] (t1) at (0.4,5.1) {}; \node[idsq] (t2) at (1.35,4.95) {}; \node[idsq] (t3) at (2.2,5.8) {}; \node[idsq] (t4) at (2.6,4.9) {};
\node[oodA] (q1) at (3.4,5.2) {}; \node[oodA] (q2) at (3.95,5.5) {}; \node[oodA] (q3) at (3.8,4.8) {}; \node[oodA] (q4) at (4.4,5.05) {};
\foreach \a/\b in {s1/s2, s2/s3, s1/t1, s1/c1, s2/c1, s3/t2, s3/c2, s3/t3, c1/t1, c1/t2, t2/c2, t3/c2, t2/t4, c2/t4, t4/q1, q1/q2, q1/q3, q2/q3, q2/q4, q3/q4, q1/q4}{\draw[edge] (\a) -- (\b);}
\draw[flow] (s1) to[bend right=12] (t1); \draw[flow] (s3) to[bend left=12] (t3); \draw[flow] (s3) to[bend right=10] (t2);
\draw[flow, orange!60!black, line width=0.5pt] (t2) to[bend right=8] (t4);
\draw[flow, orange!50!black, line width=0.35pt, densely dashed] (t4) -- (q1);
\node[lab, anchor=north west, align=left] at (4.6,6.2) {ID：支持画像に近く $u^{+}$ 大\\未知：互いに近く支持画像\\\quad から遠く $u^{+}$ 小\\参照 $\mathcal{R}$・較正 $\mathcal{C}$ も同じ節点};
\node[lab, anchor=west] at (0.0,4.5) {\tikz{\node[sup]{};}\,支持 $\mathcal{S}$\ \tikz{\node[refp]{};}\,参照 $\mathcal{R}$\ \tikz{\node[calp]{};}\,較正 $\mathcal{C}$\ \tikz{\node[idsq]{};}\,テスト（ID）\ \tikz{\node[oodA]{};}\,未知\ \tikz{\node[seed]{};}\,種};
% ================================================================ (b) BH selection on the calibrated ranks
\node[lab, anchor=west, font=\figfont\bfseries] at (0,4.0) {(b) $p^{+}$（参照画像 $\mathcal{R}$ の中での $u^{+}$ の順位）に BH 法をかけて種を選ぶ};
\draw[->, black!45] (0.3,3.2) -- (8.1,3.2);
\node[lab, anchor=north west, black!65] at (0.3,3.15) {$p^{+}$ 小（未知らしい）};
\node[lab, anchor=north east, black!65] at (8.1,3.15) {$p^{+}$ 大（ID らしい）};
\foreach \x in {0.5,0.62,0.72,0.84,0.98,1.1,1.26,1.45} {\node[seed] at (\x,3.42) {};}
\foreach \x in {1.95,2.4,3.1} {\node[oodA] at (\x,3.42) {};}
\foreach \x in {1.7,2.2,2.7,3.4,3.75,4.1,4.5,4.9,5.3,5.7,6.1,6.5,6.9,7.3,7.7} {\node[idsq] at (\x,3.41) {};}
\draw[red!70!black, line width=0.8pt, densely dashed] (1.58,3.72) -- (1.58,3.15);
\node[lab, red!70!black, anchor=west] at (1.65,3.75) {種 $\hat{Q}$：$p_{(i)}\le q\,i/(\hat{\pi}_0 m)$ を満たす最大の $i$ まで};
\node[lab, anchor=west] at (0.3,2.6) {AdaDetect の BH 法（$q{=}0.2$）：交換可能なら偽発見率 $\le q$（命題3）};
\node[lab, anchor=west] at (0.3,2.36) {枚数 $\hat{k}$ は固定しない．候補は較正画像とテスト画像のすべて（折なし）};
% ================================================================ (c) mass from the seeds and the read-out
\node[lab, anchor=west, font=\figfont\bfseries] at (0,1.6) {(c) 種の質量 $u^{-}$ を同じグラフで伝播させ，参照画像で読む};
\node[sup] (b1) at (0.55,0.85) {}; \node[sup] (b2) at (1.05,1.1) {};
\node[calp] (d1) at (0.9,0.4) {}; \node[idsq] (e1) at (0.4,0.15) {}; \node[idsq] (e2) at (1.35,0) {}; \node[idsq] (e3) at (2.2,0.85) {}; \node[idsq] (e4) at (2.6,-0.05) {};
\node[seed] (f1) at (3.4,0.25) {}; \node[seed] (f2) at (3.95,0.55) {}; \node[oodA] (f3) at (3.8,-0.15) {}; \node[oodA] (f4) at (4.4,0.1) {};
\foreach \a/\b in {b1/b2, b1/e1, b1/d1, b2/d1, b2/e3, d1/e1, d1/e2, e2/e4, e3/e4, e4/f1, f1/f2, f1/f3, f2/f3, f2/f4, f3/f4, f1/f4}{\draw[edge] (\a) -- (\b);}
\draw[nflow] (f1) to[bend right=12] (f3); \draw[nflow] (f2) to[bend left=12] (f4);
\draw[nflow, red!50!black, line width=0.35pt, densely dashed] (f1) -- (e4);
\node[lab, anchor=north west, align=left] at (4.6,1.25) {種に近い未知：$u^{-}$ 大\\\quad $\Rightarrow p^{-}$ 小\\ID：種から遠く $u^{-}$ 小\\\quad $\Rightarrow p^{-}$ 大\\種が同じクラスの\\\quad 残りを引き上げる};
\node[box, fill=yellow!15, anchor=west, align=left] (ro) at (0.3,-1.0) {$\sigma_{\mathrm{ref}}(x)=\sum_{v}\log p^{+}_{v}(x)\,p^{-}_{v}(x)$\\$s(x)=s_0(x)\exp\sigma_{\mathrm{ref}}(x)$};
\node[lab, anchor=west, align=left] at (4.3,-1.0) {$p^{\pm}$：参照画像 $\mathcal{R}$ での順位\\$p^{\mathrm{fin}}$：$\sigma_{\mathrm{ref}}$ の較正画像 $\mathcal{C}$ での順位（命題4）\\$s_0$：ベース検出器（TINS），$v$：ビュー};
\end{tikzpicture}
